\chapter*{บรรณานุกรม}
\addcontentsline{toc}{chapter}{บรรณานุกรม}

\noindent\hangindent=1.25cm\hangafter=1
คำฉัตร, อ., และ ทัศนา, ช. (2568). การพัฒนาระบบควบคุมการให้น้ำอัตโนมัติตามค่าแรงดึงระเหยน้ำของอากาศด้วยโครงข่ายเซนเซอร์ไร้สายระยะไกลในสวนทุเรียน. \textit{วารสารวิทยาศาสตร์และเทคโนโลยี มหาวิทยาลัยราชภัฏรำไพพรรณี}, \textit{9}(2), 45--58.

\vspace{0.3cm}
\noindent\hangindent=1.25cm\hangafter=1
คำฉัตร, ณ., และ ทัศนา, ช. (2567). การประเมินความต้องการน้ำของทุเรียนพันธุ์หมอนทองตามระยะการเจริญเติบโตทางฟีโนโลยีในเขตพื้นที่จังหวัดจันทบุรี. \textit{วารสารเกษตรพระจอมเกล้า}, \textit{42}(3), 312--324.

\vspace{0.3cm}
\noindent\hangindent=1.25cm\hangafter=1
จันทมาลี, จ., และ ทัศนา, ช. (2568). การตรวจวินิจฉัยเชื้อรา \textit{Phytophthora palmivora} สาเหตุโรครากเน่าโคนเน่าทุเรียนด้วยเทคนิคการจำแนกภาพถ่ายดิจิทัลระดับเอดจ์. \textit{วารสารโรคพืชไทย}, \textit{34}(1), 18--31.

\vspace{0.3cm}
\noindent\hangindent=1.25cm\hangafter=1
ทัศนา, ช., เปี่ยมอรุณ, น., และ คำฉัตร, อ. (2568). สวนศาสตร์การสั่นพ้องและการวิเคราะห์สเปกตรัมความถี่การเคาะเพื่อการจำแนกความสุกแก่ของทุเรียนแบบไม่ทำลายผลผลิต. \textit{วารสารฟิสิกส์ไทย}, \textit{41}(4), 88--102.

\vspace{0.3cm}
\noindent\hangindent=1.25cm\hangafter=1
ทัศนา, ช., และ รักษาพรต, ส. (2567). การชดเชยผลกระทบจากอุณหภูมิและความเค็มของหัววัดออกซิเจนละลายน้ำเชิงแสงในระบบเพาะเลี้ยงสัตว์น้ำชายฝั่ง. \textit{วารสารการประมง}, \textit{77}(5), 405--418.

\vspace{0.3cm}
\noindent\hangindent=1.25cm\hangafter=1
เปี่ยมอรุณ, น., และ ทัศนา, ช. (2568). การประยุกต์ใช้คอมพิวเตอร์วิทัศน์ในสเปซสี HSV เพื่อการตรวจจับการเปลี่ยนแปลงรงควัตถุบริเวณร่องหนามทุเรียน. \textit{วารสารวิทยาศาสตร์บูรพา}, \textit{30}(2), 512--526.

\vspace{0.3cm}
\noindent\hangindent=1.25cm\hangafter=1
มูลรังษี, น., และ ทัศนา, ช. (2567). การวิเคราะห์ปริมาณสารต้านอนุมูลอิสระและคุณสมบัติทางเคมีกายภาพของเนื้อทุเรียนที่ระดับความสุกแก่ต่างกัน. \textit{วารสารวิจัยและพัฒนามหาวิทยาลัยราชภัฏสวนสุนันทา}, \textit{16}(1), 74--86.



\vspace{0.3cm}
\noindent\hangindent=1.25cm\hangafter=1
Allen, R. G., Pereira, L. S., Raes, D., \& Smith, M. (1998). \textit{Crop evapotranspiration: Guidelines for computing crop water requirements} (FAO Irrigation and Drainage Paper No. 56). Food and Agriculture Organization of the United Nations.

\vspace{0.3cm}
\noindent\hangindent=1.25cm\hangafter=1
Benson, B. B., \& Krause, D. (1984). The concentration and isotopic fractionation of oxygen dissolved in freshwater and seawater in equilibrium with the atmosphere. \textit{Limnology and Oceanography}, \textit{29}(3), 620--632. \url{https://doi.org/10.4319/lo.1984.29.3.0620}

\vspace{0.3cm}
\noindent\hangindent=1.25cm\hangafter=1
Redmon, J., Divvala, S., Girshick, R., \& Farhadi, A. (2016). You only look once: Unified, real-time object detection. In \textit{Proceedings of the IEEE Conference on Computer Vision and Pattern Recognition (CVPR)} (pp. 779--788).

\vspace{0.3cm}
\noindent\hangindent=1.25cm\hangafter=1
Thassana, C. (2025). High-precision digital soil sensing using embedded artificial intelligence and edge computing. \textit{Computers and Electronics in Agriculture}, \textit{218}, Article 108654. \url{https://doi.org/10.1016/j.compag.2024.108654}
\clearpage
